\subsection{Which method? Classify first}\label{sub:odetree}
\R{read the shape, then jump}{
\begin{rcp}
\item \hd{$y'=f(x)g(y)$} (factorises) \ra\ \textbf{separation}, \S\ref{sub:separation}.
\item \hd{$y'+p(x)y=q(x)$} (linear, 1st order) \ra\ \textbf{integrating factor}, \S\ref{sub:linear1}.
\item \hd{$y'+p(x)y=q(x)y^\alpha$} \ra\ \textbf{Bernoulli}, $z=y^{1-\alpha}$, \S\ref{sub:subsode}.
\item \hd{$y'=f\!\left(\frac yx\right)$} \ra\ $u=\frac yx$, i.e.\ $y=ux$, $y'=u'x+u$ \ra\ separable.
\item \hd{$y'=f(ax+by+c)$} \ra\ $u=ax+by+c$ \ra\ separable.
\item \hd{$y^{(n)}+a_{n-1}y^{(n-1)}+\dots+a_0y=g(x)$, constant $a_i$} \ra\ characteristic polynomial (\S\ref{sub:odehom}) $+$ particular solution (\S\ref{sub:odepart}).
\item \hd{$y'=f(y)$ (autonomous)} \ra\ separable; the constant solutions are the zeros of $f$, and they decide the qualitative picture (\S\ref{sub:autonomous}).
\end{rcp}
\hd{Order of work, always:} general solution first \ra\ then insert the initial/boundary conditions \ra\ then solve for the constants.
}

\subsection{Separation of variables \de{Trennung der Variablen}}\label{sub:separation}
\R{$y'=f(x)g(y)$}{
\begin{rcp}
\item \warn \hd{First} solve $g(y_0)=0$: each root gives a \textbf{constant solution} $y\equiv y_0$. Forgetting these costs marks and they are often the ones the initial condition needs.
\item For $g(y)\ne0$: $\frac{dy}{g(y)}=f(x)\dx$, integrate both sides, one constant $C$ on the right.
\item Solve for $y$; handle $|\cdot|$ and $\pm$ by absorbing into $C$ ($e^C=:\tilde C>0$, the sign extends it to $\tilde C\in\Rl$).
\item Insert the initial condition, then state the interval of existence.
\end{rcp}
\hd{Ex.} $y'=xy$: $\frac{dy}y=x\dx$ \ra\ $\ln|y|=\frac{x^2}2+C$ \ra\ $y=\tilde Ce^{x^2/2}$ (which includes $y\equiv0$ at $\tilde C=0$).
}
\E{separable with an initial condition}{
$y'=y^2\cos x$, $y(0)=1$. Constant solution $y\equiv0$ (does not meet $y(0)=1$).\\
Otherwise $-\frac1y=\sin x+C$ \ra\ $y=\frac{-1}{\sin x+C}$; $y(0)=1$ \ra\ $C=-1$ \ra\ $\boxed{y=\frac1{1-\sin x}}$, valid near $x=0$ until $\sin x=1$.
}

\subsection{First order linear $y'+p(x)y=q(x)$}\label{sub:linear1}
\R{integrating factor (fastest, always works)}{
\begin{rcp}
\item Normal form $y'+p(x)y=q(x)$ --- the coefficient of $y'$ must be $1$.
\item $\mu(x):=e^{\int p(x)\dx}$ (no $+C$ needed).
\item $(\mu y)'=\mu q$ \ra\ $\displaystyle y=\frac1{\mu(x)}\left(\int\mu(x)q(x)\dx+C\right)$.
\item Insert the initial condition last.
\end{rcp}
\hd{``Variation of constants'' \de{Variation der Konstanten}} is the same method spelled out --- promote $C\rightsquigarrow C(x)$ in $y_h$; full recipe in \S\ref{sub:varconst}. Use the name if the exam asks for it; the integrating factor gets you there faster.
}
\E{$y'-\frac2xy=x^3$}{
$p=-\frac2x$ \ra\ $\mu=e^{-2\ln|x|}=x^{-2}$. Then $(x^{-2}y)'=x$ \ra\ $x^{-2}y=\frac{x^2}2+C$ \ra\ $\boxed{y=\frac{x^4}2+Cx^2}$.
}

\F{variation of constants \de{Variation der Konstanten} --- \textbf{1st order}, $y'+p(x)y=q(x)$}{
\label{sub:varconst}One constant, so \textbf{one} function: solve the homogeneous equation, $y_h=Ce^{-\int p}$; promote $C\rightsquigarrow C(x)$ and insert $y=C(x)e^{-\int p}$. Everything cancels except
\[C'(x)=q(x)e^{\int p},\qquad\text{so}\qquad y=e^{-\int p}\left(\int q\,e^{\int p}\dx+C\right),\]
which is exactly the integrating factor above --- use that unless the exam demands the name.
}

\subsection{Bernoulli and substitutions}\label{sub:subsode}
\R{$y'+p(x)y=q(x)y^\alpha$, $\alpha\ne0,1$}{
\begin{rcp}
\item Divide by $y^\alpha$: $y^{-\alpha}y'+py^{1-\alpha}=q$.
\item Set $z:=y^{1-\alpha}$ \ra\ $z'=(1-\alpha)y^{-\alpha}y'$.
\item \ra\ \textbf{linear} ODE $z'+(1-\alpha)pz=(1-\alpha)q$ \ra\ \S\ref{sub:linear1}.
\item Back-substitute $y=z^{1/(1-\alpha)}$; check the lost solution $y\equiv0$ separately.
\end{rcp}
\hd{Ex.} $y'+y=xy^2$ ($\alpha=2$): $z=y^{-1}$ \ra\ $z'-z=-x$ \ra\ $\mu=e^{-x}$ \ra\ $z=x+1+Ce^x$ \ra\ $y=\frac1{x+1+Ce^x}$.
}

\subsection{Constant coefficients: homogeneous part}\label{sub:odehom}
\R{$y^{(n)}+a_{n-1}y^{(n-1)}+\dots+a_0y=0$}{
\begin{rcp}
\item Ansatz $y=e^{\lambda x}$ \ra\ \hd{characteristic polynomial} $\lambda^n+a_{n-1}\lambda^{n-1}+\dots+a_0=0$.
\item \textbf{Find all roots} (guess an integer root, then polynomial division).
\item Build the fundamental system:
\begin{bul}
\item simple real root $\lambda$ \ra\ $e^{\lambda x}$
\item real root of multiplicity $m$ \ra\ $e^{\lambda x},xe^{\lambda x},\dots,x^{m-1}e^{\lambda x}$
\item complex pair $\alpha\pm i\beta$ \ra\ $e^{\alpha x}\cos\beta x,\ e^{\alpha x}\sin\beta x$ (multiplicity $m$: also times $x,\dots,x^{m-1}$)
\end{bul}
\item $y_h=C_1y_1+\dots+C_ny_n$ --- exactly $n$ free constants, no more, no fewer.
\end{rcp}
\hd{2nd order quick table} ($y''+by'+cy=0$, $D=b^2-4c$): $D>0$ \ra\ $C_1e^{\lambda_1x}+C_2e^{\lambda_2x}$; $D=0$ \ra\ $(C_1+C_2x)e^{\lambda x}$; $D<0$ \ra\ $e^{\alpha x}(C_1\cos\beta x+C_2\sin\beta x)$.
}
\subsection{Particular solution: ansatz by the right-hand side}\label{sub:odepart}
\T{ansatz table --- $m$ is the resonance multiplicity}{
{\fontsize{6.0pt}{7.1pt}\selectfont
\begin{tsTabular}{@{}>{\raggedright\arraybackslash}p{0.34\linewidth}>{\raggedright\arraybackslash}p{0.60\linewidth}@{}}
$g(x)$ & ansatz $y_p$\\
\midrule
polynomial, degree $k$ & $x^m(A_kx^k+\dots+A_0)$, full polynomial\\
$Ce^{\mu x}$ & $x^mAe^{\mu x}$\\
$C\cos\omega x$ or $C\sin\omega x$ & $x^m(A\cos\omega x+B\sin\omega x)$, \emph{both} terms\\
$P_k(x)e^{\mu x}$ & $x^me^{\mu x}(A_kx^k+\dots+A_0)$\\
$e^{\mu x}\cos\omega x$ & $x^me^{\mu x}(A\cos\omega x+B\sin\omega x)$\\
$g_1+g_2$ & solve separately and add (superposition)\\
\end{tsTabular}}
\hd{$m$ =} multiplicity of $\mu$ (resp.\ of $\mu\pm i\omega$) as a root of the characteristic polynomial. No resonance \ra\ $m=0$ \ra\ no extra $x$.
}
\R{execute the ansatz}{
\begin{rcp}
\item \textbf{Pick the shape} from the table above.
\item \warn \hd{RESONANCE CHECK:} is the exponent $\mu$ (resp.\ $\pm i\omega$) a root of the characteristic polynomial, and with what multiplicity $m$? \ra\ multiply the whole ansatz by $x^m$. Skipping this makes the ansatz collapse to $0=g$.
\item \textbf{Differentiate the ansatz}, insert into the ODE, compare coefficients, solve for $A,B,\dots$
\item \hd{General solution $=y_h+y_p$}; only now impose the initial conditions, on the \emph{sum}.
\end{rcp}
}
\E{$y''-y=e^x$ (textbook resonance)}{
\textbf{1)} $\lambda^2-1=0$ \ra\ $\lambda=\pm1$ \ra\ $y_h=C_1e^x+C_2e^{-x}$.\\
\textbf{2)} $\mu=1$ \emph{is} a simple root \ra\ resonance, $m=1$ \ra\ ansatz $y_p=Axe^x$.\\
\textbf{3)} $y_p'=A(1+x)e^x$, $y_p''=A(2+x)e^x$ \ra\ $y_p''-y_p=2Ae^x\overset!=e^x$ \ra\ $A=\frac12$.\\
\textbf{4)} $\boxed{y=C_1e^x+C_2e^{-x}+\tfrac12xe^x}$
}
\E{resonant $\cos$, and a full 2nd-order IVP}{
\hd{(a)} $y''+y=\cos x$: $\lambda=\pm i$ \ra\ $y_h=C_1\cos x+C_2\sin x$; $\pm i$ is a root \ra\ $m=1$ \ra\ $y_p=x(A\cos x+B\sin x)$ \ra\ $A=0$, $B=\frac12$ \ra\ $y=C_1\cos x+C_2\sin x+\frac x2\sin x$.\\
\hd{(b) IVP} $y''-3y'+2y=0$, $y(0)=1$, $y'(0)=0$: $\lambda=1,2$ \ra\ $y=C_1e^x+C_2e^{2x}$.
$C_1+C_2=1$, $C_1+2C_2=0$ \ra\ $C_2=-1$, $C_1=2$ \ra\ $\boxed{y=2e^x-e^{2x}}$.
}
\subsection{Conditions, existence, qualitative behaviour}\label{sub:autonomous}
\F{IVP vs BVP, Picard--Lindel\"of}{
\hd{IVP \de{Anfangswertproblem}:} all conditions at one point, $y(x_0),y'(x_0),\dots$ --- $n$ of them for order $n$.\\
\hd{BVP \de{Randwertproblem}:} conditions at two points $a,b$ --- may have $0$, exactly $1$, or infinitely many solutions.\\
\hd{Picard--Lindel\"of:} $y'=f(x,y)$, $y(x_0)=y_0$ with $f$ continuous and \emph{Lipschitz in $y$} (e.g.\ $\partial_yf$ continuous and locally bounded) \ra\ a \textbf{unique} solution on a neighbourhood of $x_0$. Higher order: same after rewriting as a system.\\
\warn Continuity of $f$ alone gives existence (Peano) but \textbf{not} uniqueness: $y'=\sqrt{|y|}$, $y(0)=0$ has both $y\equiv0$ and $y=\frac{x^2}4$.\\
Linear ODEs with continuous coefficients: unique solution on the \emph{whole} interval.
}
\R{autonomous $y'=f(y)$: sketch without solving}{
\begin{rcp}
\item Equilibria: the roots of $f(y)=0$ --- these are the constant solutions.
\item Sign of $f$ between consecutive equilibria: $f>0$ \ra\ solutions increase, $f<0$ \ra\ decrease.
\item An equilibrium $y^*$ is \hd{stable} if $f$ goes $+\to-$ there (equivalently $f'(y^*)<0$), \hd{unstable} if $-\to+$ ($f'(y^*)>0$).
\item Solutions \textbf{can never cross an equilibrium} (uniqueness), so each one is trapped between two of them and converges monotonically.
\end{rcp}
}
